\documentclass[a4paper,11ptj]{ltjsreport}

% フォントの設定
\usepackage[deluxe]{luatexja-preset}
\setmainfont{TeX Gyre Termes}
\setsansfont{TeX Gyre Heros}
\setmonofont{Inconsolatazi4}
\usepackage{amsmath}
\usepackage{unicode-math}
\setmathfont{STIX Two Math}

% 諸々のパッケージ
\usepackage{subfiles}
\usepackage{graphicx}
\usepackage{minted}

% 修論用スタイル
\usepackage{tutthesis}
\usepackage{tutabstract}

% hyperrefは最後
\usepackage[bookmarksnumbered]{hyperref}

% コードのスタイル
\setminted{
  linenos,
  breaklines,
  frame=single,
  style=friendly,
  fontsize=\small,
}

\title{タイトル}
\titleen{Title}
\idnum{211808}
\author{岩井 凜}
\authoren{Rin Iwai}
\department{情報・知能工学専攻}
\departmenten{Computer Science and Engineering}
\degree{修士（工学）}
\gyear{2026}
\supervisor{佐藤 幸紀}
\supervisoren{Yukinori Sato}
\abstractdate{2026}{12}{1}

\begin{document}
\maketitle

% 日本語アブスト
\begin{abstractja}
(要旨 1,000字程度)
\end{abstractja}

% 英語アブスト
\begin{abstracten}
(800 words)
\end{abstracten}

\pagenumbering{roman}
\setcounter{tocdepth}{3}
\tableofcontents
\listoffigures
\listoftables

\clearpage
\pagenumbering{arabic}

% 章ごとにページを羅列する
\subfile{01_introduction/introduction}

% 参考文献
\bibliographystyle{junsrt}
\bibliography{refs}

\chapter*{研究業績}
\addcontentsline{toc}{chapter}{研究業績}

\chapter*{謝辞}
\addcontentsline{toc}{chapter}{謝辞}

\appendix
\chapter{付録}

\end{document}
